\begin{document}

\newpage
\section{Unified Transfer Function Derivation}

\subsection{Voltage-to-velocity, radiation field, and electrical impedance}

For a rigid piston-like diaphragm moving with average velocity \(v_{\mathrm{avg}}(s)\), the electromechanical coupling gives
\[
v_{\mathrm{avg}}(s)=\frac{B\ell}{Z_m^{\mathrm{tot}}(s)}\,I(s)
=\frac{B\ell}{Z_m^{\mathrm{tot}}(s)}\cdot\frac{V_{\mathrm{in}}(s)}{Z_e(s)}.
\]
Thus the voltage-to-velocity transfer function is
\[
H_v(s)\triangleq\frac{v_{\mathrm{avg}}(s)}{V_{\mathrm{in}}(s)}
=\frac{B\ell}{Z_m^{\mathrm{tot}}(s)\,Z_e(s)}.
\]

The acoustic pressure at an arbitrary observation point \(\mathbf{r}\) is generated by the linear radiation operator \(\mathcal{G}_{\mathrm{rad}}(\mathbf{r},s)\):
\[
p(\mathbf{r},s)=\mathcal{G}_{\mathrm{rad}}(\mathbf{r},s)\;v_{\mathrm{avg}}(s).
\]
The operator \(\mathcal{G}_{\mathrm{rad}}(\mathbf{r},s)\) can be derived from the Rayleigh integral, a boundary-element method (BEM) solve, a Green's function, or an analytical approximation (e.g., baffled piston far-field). The voltage-to-pressure transfer function then follows as
\[
H_p(s,\mathbf{r})\triangleq\frac{p(\mathbf{r},s)}{V_{\mathrm{in}}(s)}
=\mathcal{G}_{\mathrm{rad}}(\mathbf{r},s)\,H_v(s).
\]

\subsection{General radiation operator: from surface integral to lumped model}

For a radiating surface \(S\) (the driver diaphragm) with local normal velocity \(v(\mathbf{x},s)\) at point \(\mathbf{x}\in S\), the acoustic pressure at an observation point \(\mathbf{r}\) in the Laplace domain is given by the Rayleigh integral (or its generalised form for arbitrary baffle conditions):
\[
\boxed{
p(\mathbf{r},s) = \frac{s\rho_0}{2\pi} \iint_S v(\mathbf{x},s) \, \frac{e^{-k R(\mathbf{r},\mathbf{x})}}{R(\mathbf{r},\mathbf{x})} \, dS(\mathbf{x})
},
\]
where \(k = s/c\) and \(R(\mathbf{r},\mathbf{x}) = |\mathbf{r} - \mathbf{x}|\). This is the most general linear relation between the surface velocity distribution and the field pressure.

The operator \(\mathcal{G}_{\mathrm{rad}}(\mathbf{r},s)\) is therefore an integral operator:
\[
p(\mathbf{r},s) = \int_S G(\mathbf{r},\mathbf{x},s)\, v(\mathbf{x},s) \, dS(\mathbf{x}),
\]
with kernel \(G(\mathbf{r},\mathbf{x},s) = \frac{s\rho_0}{2\pi} \frac{e^{-kR}}{R}\) (for an infinite baffle; modify for other boundary conditions).

\subsubsection{Collapsed piston approximation}

If the diaphragm moves as a rigid piston with uniform velocity \(v_{\mathrm{avg}}(s) = \frac{1}{S_d}\iint_S v(\mathbf{x},s)\,dS\) and the observation point is in the far field (or the surface is small compared to wavelength), the integral simplifies:
\[
p(\mathbf{r},s) \approx \frac{s\rho_0 S_d}{2\pi r}\, e^{-sr/c}\, v_{\mathrm{avg}}(s) \quad \text{(far-field, baffled)}.
\]
This yields the familiar scalar operator \(\mathcal{G}_{\mathrm{rad}}(\mathbf{r},s) \approx \frac{\rho_0 S_d s}{\Omega r} e^{-sr/c}\) with \(\Omega = 2\pi\) for half‑space.

\subsubsection{When the collapsed model fails}

The uniform‑velocity assumption breaks down when:
\begin{itemize}
  \item The diaphragm exhibits **cone breakup** (flexural modes) — typical above a few kilohertz.
  \item The velocity distribution is intentionally non‑uniform (e.g., dome tweeters with surround elasticity).
  \item Near‑field or high‑precision modelling is required (e.g., microphone calibration, array design).
\end{itemize}
In such cases, the full integral must be used, which can be discretised via BEM or ring‑element methods.

\subsubsection{Consistency with mechanical radiation impedance}

The mechanical radiation impedance \(Z_{\mathrm{rad}}(s)\) is defined from the total force exerted by the fluid on the surface:
\[
F_{\mathrm{rad}}(s) = \iint_S p_{\mathrm{near}}(\mathbf{x},s) \, dS(\mathbf{x}),
\]
where \(p_{\mathrm{near}}\) is the pressure on the surface. For a rigid piston with uniform velocity, this reduces to the scalar relation \(Z_{\mathrm{rad}}(s) = S_d \, Z_{\mathrm{rad}}^{\mathrm{ac}}(s)\). For general velocity distributions, it becomes an impedance matrix coupling different parts of the surface.

In the modular decomposition \(Z_m^{\mathrm{tot}} = Z_{\mathrm{driver}} + Z_{\mathrm{rad}} + Z_{\mathrm{enclosure}}\), the term \(Z_{\mathrm{rad}}(s)\) is now understood as the **effective scalar radiation impedance** seen by the average velocity. This effective impedance can be obtained by projecting the full impedance matrix onto the assumed mode shape (usually the rigid‑body mode). If breakup modes are important, the model must be extended to multiple degrees of freedom.

\subsubsection{Practical implementation for rigid pistons}

For most loudspeaker design up to the breakup frequency, the rigid‑piston assumption is sufficient. In that case:
\begin{itemize}
  \item Use the full Rayleigh integral (numerically integrated over the cone) to compute both \(Z_{\mathrm{rad}}(s)\) and \(\mathcal{G}_{\mathrm{rad}}(\mathbf{r},s)\) consistently.
  \item The average velocity is then the only mechanical state variable, and the collapsed formulation is exact (within the rigid assumption).
\end{itemize}
For a circular piston in an infinite baffle, the on‑axis pressure and the radiation impedance have closed forms (Bessel/Hankel); for arbitrary shapes or baffles, use numerical quadrature or BEM.

\subsection{Total mechanical impedance — modular decomposition}

The total mechanical impedance seen by the motor is the sum of three physically independent contributions:
\[
\boxed{Z_m^{\mathrm{tot}}(s)=Z_{\mathrm{driver}}(s)+Z_{\mathrm{rad}}(s)+Z_{\mathrm{enclosure}}(s)}.
\]

\noindent
\textbf{1. Pure driver impedance} (intrinsic mechanical motion):
\[
\boxed{Z_{\mathrm{driver}}(s)=R_{\mathrm{ms}}+sM_{\mathrm{ms}}+\frac{1}{sC_{\mathrm{ms}}}}.
\]
This contains only the driver's own mechanical resistance, moving mass, and suspension compliance — the parameters measured with no external acoustic load (except free air, which already includes low-frequency added mass; see note below).

\noindent
\textbf{2. Full radiation impedance} (mechanical side):
\[
\boxed{Z_{\mathrm{rad}}(s)=R_{\mathrm{rad}}(s)+sM_{\mathrm{ma}}(s)}.
\]
Here
\begin{itemize}
  \item \(M_{\mathrm{ma}}(s)\) is the frequency-dependent added mass (reactive near-field energy storage).
  \item \(R_{\mathrm{rad}}(s)\) is the frequency-dependent radiation resistance (real power radiated away).
\end{itemize}
Both arise from the same acoustic problem and are obtained together — e.g., as the real and imaginary parts of the self‑loading computed via BEM or the Rayleigh integral.

\noindent
\textbf{3. Enclosure impedance} (rear loading and internal losses):
\[
\boxed{Z_{\mathrm{enclosure}}(s)=\text{rear acoustic loading reflected to the mechanical side}}.
\]
This accounts for any structure that loads the diaphragm from the rear: sealed-box compliance, vented-box Helmholtz resonator, stuffing resistance, passive radiator, or a rear horn. For an open baffle (no enclosure), \(Z_{\mathrm{enclosure}}(s)=0\).

The total mechanical impedance explicitly becomes
\[
Z_m^{\mathrm{tot}}(s)=
s\bigl(M_{\mathrm{ms}}+M_{\mathrm{ma}}(s)\bigr)
+R_{\mathrm{ms}}
+\frac{1}{sC_{\mathrm{ms}}}
+R_{\mathrm{rad}}(s)
+Z_{\mathrm{enclosure}}(s).
\]

\subsection{Electrical impedance}

The electrical impedance at the voice-coil terminals includes the reflected motional term:
\[
\boxed{Z_e(s)=R_e+sL_e+\frac{(B\ell)^2}{Z_m^{\mathrm{tot}}(s)}}.
\]
Substituting the explicit mechanical impedance yields
\[
\boxed{
Z_e(s)=R_e+sL_e+
\frac{(B\ell)^2}{
s\bigl(M_{\mathrm{ms}}+M_{\mathrm{ma}}(s)\bigr)
+R_{\mathrm{ms}}
+\dfrac{1}{sC_{\mathrm{ms}}}
+R_{\mathrm{rad}}(s)
+Z_{\mathrm{enclosure}}(s)
}.
}
\]

\subsection{General voltage-to-pressure transfer function}

Combining the above gives the most general linear transfer function:
\[
\boxed{
H_p(s,\mathbf{r})=
\frac{p(\mathbf{r},s)}{V_{\mathrm{in}}(s)}
=
\frac{B\ell\;\mathcal{G}_{\mathrm{rad}}(\mathbf{r},s)}
{(B\ell)^2+\bigl(R_e+sL_e\bigr)\,Z_m^{\mathrm{tot}}(s)}
}.
\]
Equivalently, with the explicit mechanical impedance:
\[
\boxed{
H_p(s,\mathbf{r})=
\frac{B\ell\;\mathcal{G}_{\mathrm{rad}}(\mathbf{r},s)}
{(B\ell)^2+\bigl(R_e+sL_e\bigr)
\left[
s\bigl(M_{\mathrm{ms}}+M_{\mathrm{ma}}(s)\bigr)
+R_{\mathrm{ms}}
+\dfrac{1}{sC_{\mathrm{ms}}}
+R_{\mathrm{rad}}(s)
+Z_{\mathrm{enclosure}}(s)
\right]
}.
}
\]

\subsection{Relation to classical Thiele–Small models}

Classical T/S is recovered by:
\begin{itemize}
  \item Using the low-frequency, directivity-free far-field approximation
    \[
    \mathcal{G}_{\mathrm{rad}}(\mathbf{r},s)\approx \frac{\rho_0 S_d s}{\Omega r}\,e^{-sr/c},
    \]
    where \(\Omega=2\pi\) for half-space (infinite baffle), \(4\pi\) for free space, and other values only as geometric low-frequency factors.
  \item Neglecting the frequency dependence of \(M_{\mathrm{ma}}(s)\) (set to constant low-frequency added mass \(\frac{8}{3}\rho_0 a^3\)).
  \item Ignoring \(R_{\mathrm{rad}}(s)\) entirely (justified at low frequencies where \(R_{\mathrm{rad}}\ll R_{\mathrm{ms}}\)).
  \item Selecting the appropriate \(Z_{\mathrm{enclosure}}(s)\) for sealed or vented boxes.
\end{itemize}
For many modern applications (wideband drivers, open baffles, horns, or high-frequency simulation), the full modular form above is necessary.

\subsection{Practical enclosure and loading forms}

All enclosure impedances are expressed on the mechanical side (force per velocity). Conversion from acoustic impedance uses the factor \(S_d^2\) where appropriate.

\paragraph{Sealed enclosure (undamped):}
\[
\boxed{Z_{\mathrm{enclosure}}^{\text{sealed}}(s)=\frac{1}{s\,C_{\mathrm{ab}}\,S_d^2}},\qquad
C_{\mathrm{ab}}=\frac{V_b}{\rho_0 c^2}.
\]

\paragraph{Sealed enclosure with damping (stuffing):}
A simple model places a stuffing resistance in parallel with the box compliance:
\[
\boxed{Z_{\mathrm{enclosure}}^{\text{sealed+damped}}(s)
=\left(\frac{1}{s\,C_{\mathrm{ab}}\,S_d^2}\right)\parallel R_{\mathrm{stuff}}(s)}.
\]
More accurate models use frequency-dependent \(R_{\mathrm{stuff}}(s)\) or a series resistance; the parallel form is a convenient first approximation.

\paragraph{Ported (vented) enclosure:}
\[
\boxed{Z_{\mathrm{enclosure}}^{\text{ported}}(s)
=S_d^2\left(
\frac{1}{sC_{\mathrm{ab}}}
\parallel
\bigl(sM_{\mathrm{ap}}+R_{\mathrm{ap}}\bigr)
\right)},
\]
where \(M_{\mathrm{ap}}\) is the acoustic mass of the port (including end corrections) and \(R_{\mathrm{ap}}\) accounts for viscous and radiation losses at the port.

\paragraph{Passive radiator:}
A passive radiator adds an additional mass-spring-damper branch coupled to the box compliance. Its mechanical impedance (referred to the driver cone) is
\[
Z_{\mathrm{pr}}(s)=\frac{S_d^2}{S_{\mathrm{pr}}^2}\left(sM_{\mathrm{pr}}+\frac{1}{sC_{\mathrm{pr}}}+R_{\mathrm{pr}}\right),
\]
and it appears in parallel with the box compliance inside \(Z_{\mathrm{enclosure}}(s)\).

\paragraph{Horn loading:}
A front-loaded horn is best treated as part of the radiation operator \(\mathcal{G}_{\mathrm{rad}}(\mathbf{r},s)\) and the radiation impedance \(Z_{\mathrm{rad}}(s)\). For a rear-loaded horn, it contributes to \(Z_{\mathrm{enclosure}}(s)\) via a transmission-line or transfer-matrix model. In either case, the modular separation remains valid.

\subsection{Key physical relationships}

\begin{align*}
\text{Velocity to pressure:} &\quad p(\mathbf{r},s)=\mathcal{G}_{\mathrm{rad}}(\mathbf{r},s)\,v(s),\\
\text{Displacement to pressure:} &\quad p(\mathbf{r},s)=\mathcal{G}_{\mathrm{rad}}(\mathbf{r},s)\,s\,x(s),\\
\text{Acceleration to pressure:} &\quad p(\mathbf{r},s)=\mathcal{G}_{\mathrm{rad}}(\mathbf{r},s)\,\frac{1}{s}\,a(s).
\end{align*}
Thus the three forms are consistent; only the kinematic variable differs.

\subsection{Modeling and implementation notes}

\begin{itemize}
  \item \textbf{Consistency} between \(\mathcal{G}_{\mathrm{rad}}(\mathbf{r},s)\) and \(Z_{\mathrm{rad}}(s)\): They must derive from the same acoustic model (e.g., the same BEM solve) to maintain reciprocal coupling between mechanical loading and sound generation.
  \item \textbf{Propagation delay} is naturally included in \(\mathcal{G}_{\mathrm{rad}}(\mathbf{r},s)\) when using the retarded Green's function (e.g., \(e^{-kR}/R\) in the Laplace domain). For far-field algebraic approximations, retain the factor \(e^{-sr/c}\) explicitly.
  \item \textbf{Non-rigid diaphragms} (cone breakup, flexural modes) invalidate the scalar average velocity assumption. For such cases, a distributed surface velocity model or a multi-degree-of-freedom modal expansion is required, but the modular structure can be extended by replacing the scalar \(v_{\mathrm{avg}}(s)\) with a vector of modal velocities.
  \item \textbf{Numerical implementation:} Compute \(Z_{\mathrm{rad}}(s)\) via BEM or the Rayleigh integral on the desired frequency grid. Store both real and imaginary parts as two vectors. Use rational approximation (e.g., vector fitting) if a state-space model is needed for time-domain simulation.
  \item \textbf{Low-frequency limits:} When \(ka\ll1\), \(M_{\mathrm{ma}}\) approaches a constant \(\frac{8}{3}\rho_0 a^3\) (for a baffled piston) and \(R_{\mathrm{rad}}\) is proportional to \(\omega^2\). In that regime, the classical T/S model works well.
\end{itemize}

\end{document}